\documentclass{article}
\usepackage{graphicx,xy,amsmath,amssymb,amsthm,physics,mathtools,tcolorbox,hyperref}
\usepackage{xepersian}
\settextfont{XB Niloofar}
\renewcommand{\contentsname}{هفته های تدریس}
\title{	
	پیش گزارش دوم درس آزمایشگاه اپتیک - دکتر مهدوی
	\\
	\small
	موضوع آزمايش: اندازه گیری طول موج با استفاده از منحنی پاشندگی
	نوری
}
\author{
حسین محمدی 
\\
۹۶۱۰۱۰۳۵
}
\begin{document}
\maketitle
\section{هدف آزمایش}
هدف غایی این آزمایش، بدست آوردن طول موجِ نوری است که در آزمایشگاه به وسیله لامپ هلیوم یا کادمیم در اختیار ما قرار می گیرد. اما در حین انجام این آزمایش با توری پراش که از شگفت انگیزترین و زیباترین وسیله در آزمایشگاه اپتیک است آشنا می شویم و معادله پراش را یاد می گیریم.

\noindent
در حقیقت سعی می کنیم با بررسی الگوی پراشی که ایجاد شده است، طول موج های مختلف موجود در طیف هلیوم یا کادمیم را بررسی کنیم.
\section{چگونگی انجام آزمایش}
باز هم نیاز به طیف سنج داریم که از قبل طرز کار آن را بلدیم و وسیله ای که اضافه بر آزمایش قبل داریم، توری پراش است.(و لامپ کادمیم.)

ابتدا بایستی توری پراش رو طوری روی صفحه حامل طیف سنج نصب کنیم که باریکه نوری که از دوربین موازی ساز خارج می شود،‌دقیقا عمود بر سطح توری پراش بتابد و این کار با نشان کردن یک خط خاص(مثلا خط قرمز از مرتبه سوم پراش) انجام می دهیم و زاوایایی را که از مشاهده این خط طیفی در دو طرف توری پراش بدست می آید را از هم کم می کنیم و حامل را به اندازه نصف اختلاف آن زوایا می چرخانیم، پس از آن نور خروجی از دروبین موازی ساز دقیقا بر صفحه پراش عمود می شود.

برای انجام آزمایش اصلی، لامپ هلیوم را پشت دریچه موازی ساز قرار می دهیم و سپس اندازه ی زاویه پراش  مرتبه اول را برای طیف اتم هلیوم(یعنی تک تک خط های موجود در طیف) انجام می دهیم و مقادیرش را یادداشت می کنیم. همین اندازه گیری ها را برای مرتبه دوم و سوم پراش و همچنین با تغییر لامپ به لامپ کادمیم انجام می دهیم و داده های مربوط به زاویه پراش را می خوانیم.

با داشتن زاویه پراش با جایگذاری در معادله توری پراش یعنی
\[
a (\sin (\alpha) - \sin (i) ) = k \lambda
\]
با داشتن زاویه پراش و بازتاب و داشتن فاصله بین شیار متوالی روی توری پراش و داشتن مرتبه پراش، یافتن طول موج کار آسانی است.
\section{مهم ترین کاربرد توری پراش}
توری پراش کارش تجزیه نور است و مهمترین کاربرد آن در طیف شناسی است، دقیقا همان کاری که اینجا کردیم، یعنی طیف اتم هلیوم و کادمیم را شناختیم، به طور مشابه می توان طیف نشری هر ماده ای را با کمک این روش شناسایی کرد.

حتی برای شناخت ساختار اجرام آسمانی از این روش استفاده می شود. عنصر هلیوم برای اولین بار مطابق همین روش کشف شد.
\section{علت عدم انطباق رنگهای مختلف مرتبه های غیر صفر توری پراش}
برای پاسخ دادن به این سوال نیاز به دانستن چیزهای جدید داریم؛ 

اگر بخواهم به صورت سر دستی جواب بدهم، پاسخ این است که محل تداخل های سازنده روی پرده به طول موج بستگی دارد (از آزمایش دوشکاف یانگ می دانیم که فاصله n امین نوار روشن از مرکز پرده برابر با 
$y_n = n \lambda \frac{d}{D}$
است که 
\lr{d}
فاصله دو شکاف از یکدیگر و 
\lr{D}
فاصله دو شکاف از پرده است.
)
پس برای یک موج مرکب، بدیهی است که فاصله ی نورهای با طول موج مختلف باید از مرکز پرده متفاوت باشد و این دقیقا یعنی که برای مرتبه غیر صفر پراش، طول موجهای مختلف از یکدیگر جدا می شوند.

\noindent\\
اما جواب دقیقتر را که با جستجو در اینترنت به آن پی بردم، شیوه کارکرد توری های پراش با «اصل هویگنس» است.

مطابق این اصل، هر نقطه ای چشمه موج را می توان به عنوان یک منبع موج کروی در نظر گرفت و در هنگام بازتاب یا در هنگام شکست نور می توان از این تناظر استفاده کرد، توجه کنید که جبهه موج اصلی «پوش»  منحنی های فاز ثابت این موج های کروی است. (برای این که منظورم را خوب متوجه شوید، یک انیمیشن خیلی خوب در اینترنت پیدا کردم که لینکش را اینجا قرار می دهم و به کمک آن می توان شکست نور در مرز محیط و بازتابش نور را به خوبی مشاهده کرد.
\href{https://www.walter-fendt.de/html5/phen/refractionhuygens_en.htm}{مشاهده لینک}
)

\noindent\\
در اینجا هم می توان مطابق اصل هوینگس این تحلیل تصویری را برای بازتاب از سطح داشت(شکل 
\ref{Fig1}
):
\begin{figure}
	\centering
	\includegraphics[width=\linewidth]{1.jpg}
	\caption{چگونگی بازتاب مطابق اصل هویگنس}
	\label{Fig1}
\end{figure}

حالا مطابق شکل
\ref{Fig1}
فقط یک دسته پوش وجود ندارد،‌بلکه می توان پوشهای شکل
\ref{Fig2}
را هم در نظر گرفت و این دقیقا همان چیزی است که ما از آن به عنوان پراش های مرتبه دوم و ... تعبیر می کنیم.(توجه کنید که رنگ متفاوت، نشانگر رنگ موج نیست بلکه برای تفکیک بهتر از رنگ های متفاوت استفاده شده است.)

\begin{figure}
	\centering
	\includegraphics[width=\linewidth]{2.jpg}
	\caption{مرتبه های مختلف پراش با کمک اصل هویگنس}
	\label{Fig2}
\end{figure}

\noindent
حالا توجه داشته باشید فاصله ی بین دایره های متوالی در تصویر، همان طول موج است و وقتی که طول موج متفاوت باشد،زاویه بازتاب شدن و پراش از مرتبه های بعد از صفر، نسبت به طول موج های دیگر کمی منحرف می شود و این همان پاشندگی نور را به ما می دهد.(توجیه پراش به کمک اصل هویگنس)















\end{document}